% figures/threshold_radar.tex
% 临界阈值蛛网图：四轴雷达图（发射成本/GPU能效/散热密度/硬件寿命）
% 使用 TikZ 手工绘制四边形雷达图

\begin{tikzpicture}[scale=1.0]

% ===== 轴定义 =====
% 四个轴：右(0°)、上(90°)、左(180°)、下(270°)
% 中心点 (0,0)，最大半径 3.5cm

\def\rmax{3.5}

% ===== 网格线（0.25, 0.5, 0.75, 1.0） =====
\foreach \r in {0.875,1.75,2.625,3.5}{
    \draw[gray!30, line width=0.3pt]
        (\r,0) -- (0,\r) -- (-\r,0) -- (0,-\r) -- cycle;
}

% ===== 四轴 =====
% 轴 1: 右 — GPU 能效 (TFLOPS/W)
% 轴 2: 上 — 散热密度 (W/m²)
% 轴 3: 左 — 发射成本 ($/kg, 反向)
% 轴 4: 下 — 硬件寿命 (年)
\draw[->, >=stealth, graybase, line width=0.8pt] (0,0) -- ( \rmax,0);
\draw[->, >=stealth, graybase, line width=0.8pt] (0,0) -- (0, \rmax);
\draw[->, >=stealth, graybase, line width=0.8pt] (0,0) -- (-\rmax,0);
\draw[->, >=stealth, graybase, line width=0.8pt] (0,0) -- (0,-\rmax);

% ===== 轴标签 =====
\node[primary, font=\bfseries\footnotesize, anchor=west]
    at (\rmax+0.3,0) {\textbf{GPU能效}};
\node[primary, font=\scriptsize, anchor=west]
    at (\rmax+0.3,-0.35) {TFLOPS/W};

\node[primary, font=\bfseries\footnotesize, anchor=south]
    at (0,\rmax+0.3) {\textbf{散热密度}};
\node[primary, font=\scriptsize, anchor=south]
    at (0,\rmax+0.05) {W/m\textsuperscript{2}};

\node[primary, font=\bfseries\footnotesize, anchor=east]
    at (-\rmax-0.3,0) {\textbf{发射成本}};
\node[primary, font=\scriptsize, anchor=east]
    at (-\rmax-0.3,-0.35) {\$/kg (反向)};

\node[primary, font=\bfseries\footnotesize, anchor=north]
    at (0,-\rmax-0.3) {\textbf{硬件寿命}};
\node[primary, font=\scriptsize, anchor=north]
    at (0,-\rmax-0.05) {年};

% ===== 刻度标记 =====
% GPU能效: 0, 25, 50, 75, 100, 125, 150
\foreach \v/\l in {0.583/25, 1.167/50, 1.75/75, 2.333/100, 2.917/125, 3.5/150}{
    \draw[gray!40] (\v,0.08) -- (\v,-0.08);
    \node[font=\tiny, graybase, anchor=north] at (\v,-0.1) {\l};
}

% 散热密度: 0, 500, 1000, 1500, 2000, 2500, 3000
\foreach \v/\l in {0.583/500, 1.167/1000, 1.75/1500, 2.333/2000, 2.917/2500, 3.5/3000}{
    \draw[gray!40] (0.08,\v) -- (-0.08,\v);
    \node[font=\tiny, graybase, anchor=east] at (-0.12,\v) {\l};
}

% 发射成本 (反向): 0, 50, 100, 150, 200, 250, 300
\foreach \v/\l in {0.583/50, 1.167/100, 1.75/150, 2.333/200, 2.917/250, 3.5/300}{
    \draw[gray!40] (-\v,0.08) -- (-\v,-0.08);
    \node[font=\tiny, graybase, anchor=north] at (-\v,-0.1) {\l};
}

% 硬件寿命: 0, 2, 4, 6, 8, 10, 12
\foreach \v/\l in {0.583/2, 1.167/4, 1.75/6, 2.333/8, 2.917/10, 3.5/12}{
    \draw[gray!40] (0.08,-\v) -- (-0.08,-\v);
    \node[font=\tiny, graybase, anchor=east] at (-0.12,-\v) {\l};
}

% ===== 红色虚线多边形: 2026 基准 =====
% 发射 $200 → 左轴: 200/300 * 3.5 = 2.333
% GPU能效 7.5 → 右轴: 7.5/150 * 3.5 = 0.175
% 散热 1400 → 上轴: 1400/3000 * 3.5 = 1.633
% 寿命 5 → 下轴: 5/12 * 3.5 = 1.458
\fill[highlight, opacity=0.15]
    (0.175,0) -- (0,1.633) -- (-2.333,0) -- (0,-1.458) -- cycle;
\draw[highlight, line width=1.2pt, dashed]
    (0.175,0) -- (0,1.633) -- (-2.333,0) -- (0,-1.458) -- cycle;
\node[highlight, font=\footnotesize\bfseries, anchor=south west]
    at (0.2,0.2) {\textbf{2026 基准}};

% 数据点标注
\fill[highlight] (0.175,0) circle (1.8pt);
\fill[highlight] (0,1.633) circle (1.8pt);
\fill[highlight] (-2.333,0) circle (1.8pt);
\fill[highlight] (0,-1.458) circle (1.8pt);

% ===== 蓝色实线多边形: 翻转条件 =====
% 发射 <$100 → 100/300 * 3.5 = 1.167
% GPU能效 >100 → 100/150 * 3.5 = 2.333
% 散热 >2000 → 2000/3000 * 3.5 = 2.333
% 寿命 >7 → 7/12 * 3.5 = 2.042
\fill[primary, opacity=0.25]
    (2.333,0) -- (0,2.333) -- (-1.167,0) -- (0,-2.042) -- cycle;
\draw[primary, line width=1.5pt]
    (2.333,0) -- (0,2.333) -- (-1.167,0) -- (0,-2.042) -- cycle;
\node[primary, font=\footnotesize\bfseries, anchor=south east]
    at (1.5,1.5) {\textbf{翻转条件}};

% 数据点标注
\fill[primary] (2.333,0) circle (2pt);
\fill[primary] (0,2.333) circle (2pt);
\fill[primary] (-1.167,0) circle (2pt);
\fill[primary] (0,-2.042) circle (2pt);

% ===== 面积差标注 =====
\node[graybase, font=\tiny\itshape, align=center]
    at (0.8,0.6) {\textbf{面积差}\\\textbf{$=$ 不可行区间}};

% ===== 图例 =====
\node[font=\footnotesize, anchor=north west]
    at (-3.5,3.5) {
    \begin{tabular}{ll}
    \tikz{\draw[highlight, dashed, line width=1.2pt] (0,0) -- (0.5,0);} & 2026 基准 \\
    \tikz{\draw[primary, line width=1.5pt] (0,0) -- (0.5,0);} & 翻转条件 \\
    \end{tabular}
};

\end{tikzpicture}
